\section{Level scores and how we compute them}
\label{sec:estimators}
% Everything in this section is condensed from docs/paper/score-spec.md (`audit-v1` scorer
% spec, draft for sign-off, 18 Oct 2026) and implemented in src/apexzero/audit/estimators.py.
% Section numbers "spec §x" refer to that file.

An audit of a score is only meaningful if the score is the one practitioners run. We
therefore port every score from pinned reference code, keep its behaviour wherever our
setting allows, and list every departure. The references are JaxUED at commit
\code{0f8f128} for PVL, MaxMC and the \PLRp{} buffer, minimax at \code{2ae9e04} for L1 value
loss and ensemble disagreement, the MNA authors' repository at \code{12340c7}, and the SFL
authors' repository at \code{8b14e2e}~\citep{coward2024jaxued,minimax2023code,mna2026code,sfl2024code}.
% src: docs/paper/score-spec.md section 0.1 (pinned sources)

\subsection{Notation and the scoring advantage}
\label{sec:est-notation}

One episode $e$ is one window, running over steps $t_0(e), \dots, t_{\text{end}}(e)$ with
$T_e$ steps including the terminal transition, as in JaxUED. $r_t$ is the raw reward of
\cref{eq:reward}. Our PPO trains its critic on rewards divided by a running return
standard deviation $\sigma$, so we rescale every value to raw units, $V_t := \sigma \hat
V_t$, and likewise the five ensemble heads $V^k_t$ (\dev{6}). $\mathrm{term}_t$ marks a
crash and $\mathrm{trunc}_t$ the window end. $V^f_t$ is the value of the pre-reset final
observation at a truncation, and $\kappa_t = 1$ marks the last step of a 128-step training
chunk that does not end the episode. In the fresh audit, window starts are given a random
phase $\varphi_e \sim U\{0, \dots, 127\}$ against chunk boundaries, as they have in
training. We write $G_e = \sum_{t \in e} r_t$, $G^\gamma_e = \sum_{\tau=0}^{T_e-1}
\gamma^\tau r_{t_0+\tau}$ (no bootstrap) and $\bar V_e = T_e^{-1}\sum_{t\in e} V_t$.
% src: docs/paper/score-spec.md section 0.2

The reference codebases score levels with the PPO's own GAE tensor. Ours is in scaled
units and bootstraps truncations, so we run a separate scoring pass in raw units
(\dev{6}):
\begin{equation}
  \mathrm{nxt}_t =
  \begin{cases}
    0 & \text{if } \mathrm{term}_t,\\
    0 \ \text{(mode F) or}\ V^f_t \ \text{(mode B)} & \text{if } \mathrm{trunc}_t,\\
    V_{t+1} & \text{otherwise},
  \end{cases}
  \qquad
  \begin{aligned}
  \delta_t &= r_t + \gamma\,\mathrm{nxt}_t - V_t,\\
  A_t &= \delta_t + \gamma\lambda\,(1-d_t)(1-\kappa_t)\,A_{t+1},
  \end{aligned}
  \label{eq:gae}
\end{equation}
with $d_t = \mathrm{term}_t \lor \mathrm{trunc}_t$, $\gamma = 0.995$ and the PPO's
$\lambda = 0.95$. At a chunk end ($\kappa_t = 1$) the $\lambda$-sum restarts and $V_{t+1}$,
computed with the parameters that produced $V_t$, bootstraps it.
% src: docs/paper/score-spec.md section 0.3; src/apexzero/audit/estimators.py (scoring_gae)
Mode F treats the window end as terminal, as JaxUED, minimax and the MNA code all do with
their time limits; it is the headline. Mode B bootstraps, as our PPO does, and is stored
as a variant.

\paragraph{An expected artefact of mode F} The critic has no time-in-window input, so at
every window end $\delta \approx r - V$, roughly $-60$ in raw units at racing speed. This
pushes $A_t$ strongly negative over the last $1/(1-\gamma\lambda) \approx 18$ steps, adds
about $|V_{\text{end}}| \cdot 18 / T_e$ to L1, and biases the mean advantage negative, so
$\mathrm{PVL}/\mathrm{L1} < 0.5$ even with an accurate critic. The prediction that this
ratio tends to $0.5$ as critic error shrinks can therefore only hold in mode B. We log
both modes and the tail term, the summed $|A_t|$ over the last 64 steps of each completed
window.
% src: docs/paper/score-spec.md section 0.3 ("Notes on mode F") and section 6 ("E2 wording")

\subsection{The six headline estimators}
\label{sec:est-six}

\Cref{tab:estimators} gives each score as we compute it per level from a set $\mathcal E$
of $E$ completed windows. Scores are episode-weighted means, as in the reference
accumulator. We store every score for $\lambda \in \{0.9, 0.95, 0.98, 1\}$, in mode B, and
in a discount-weighted form $\sum_\tau \gamma^\tau m_\tau / \sum_\tau \gamma^\tau$; the
appendix variants are $1-p$, return variance and step-normalised absolute learning
progress.
% src: docs/paper/score-spec.md sections 1.1-1.8; docs/plan/two-track-plan.md section 2.2
%      ("Six headline estimators", "Stored variants")
\TODO{the three appendix variants are named in the plan but not yet specified or
implemented}

\begin{table}[t]
\centering
\small
\caption{The six headline scores (seven rows: MaxMC has two variants). ``Reference'' is
the pinned function we port. All are larger for levels the score deems more worth
replaying.}
\label{tab:estimators}
% src: docs/paper/score-spec.md sections 1.1-1.7; src/apexzero/audit/estimators.py (level_scores)
\renewcommand{\arraystretch}{1.25}
\begin{tabularx}{\linewidth}{@{}>{\raggedright\arraybackslash}p{2.2cm}>{\raggedright\arraybackslash}X>{\raggedright\arraybackslash}p{2.9cm}@{}}
\toprule
Score & Per-level value & Reference \\
\midrule
L1 value loss & $\frac{1}{E}\sum_e \frac{1}{T_e}\sum_{t\in e} |A_t|$ &
  minimax \fn{compute_l1_value_loss} \\
PVL & $\frac{1}{E}\sum_e \frac{1}{T_e}\sum_{t\in e} \max(A_t, 0)$ &
  JaxUED \fn{positive_value_loss} \\
MaxMC (as implemented) & $\Rmax - \frac{1}{E}\sum_e \bar V_e$, \ $\Rmax = \max\!\big(0,
  \max_e G_e\big)$: undiscounted and floored at zero &
  JaxUED \fn{max_mc}, \fn{compute_max_returns} \\
MaxMC (unit-consistent) & $\frac{1}{E}\sum_e \big(\Rgmax - V_{t_0(e)}\big)$, \
  $\Rgmax = \max_e G^\gamma_e$: discounted, no floor, no time average & ours (\dev{22}) \\
Learnability & $\hat p(1-\hat p)$, $\hat p$ the mean of a per-window success indicator:
  S1 = the window completes; S2 = S1 and distance $\ge 0.97\,d_{\text{QSS}}(\theta)$ &
  SFL \fn{_calc_outcomes_by_agent} \\
MNA & $\hat C \cdot \sum_e\sum_{t\in e}\hat G_t \,/\, \sum_e T_e$ with the closed-form
  clipped advantage $\hat G$ below; $\hat C = \ind[\text{any window satisfied S1}]$ &
  MNA \fn{compute_clipped_gae}, \fn{negative_value_loss_solved_sum} \\
Disagreement & $\frac{1}{E}\sum_e\frac{1}{T_e}\sum_{t\in e}\std_k V^k_t$ over the five
  heads (population std; main critic excluded) & minimax
  \fn{compute_value_disagreement} \\
\bottomrule
\end{tabularx}
\end{table}

\paragraph{MaxMC's units} As implemented, MaxMC subtracts a discounted value (at most
$1/(1-\gamma) = 200$ steps of reward) from an undiscounted return over up to \num{1200}
steps, and floors the maximum at zero because the reference accumulator starts at zero.
We keep this because it is what runs; the unit-consistent variant is our addition and is
labelled as such.
% src: docs/paper/score-spec.md sections 1.3, 1.4; docs/paper/sources-jaxued.md section 1.1

\paragraph{Learnability under dense reward} A literal port of SFL's success count would
count every step with positive reward and give $p > 1$, so success is a per-window
indicator (\dev{9}). S1, window completion, saturates as the agent becomes competent;
S2 also asks for at least 97\% of the distance the quasi-steady-state oracle covers from
$x_0$ in 60\,s, and does not saturate. Because the tie fraction of $p(1-p)$ follows from
how success is defined, we report it only descriptively.
% src: docs/paper/score-spec.md section 1.5; docs/plan/two-track-plan.md section 2.5
\TODO{define $d_{\text{QSS}}$ precisely (spec section 6, open item)}

\paragraph{MNA in closed form} Let $A^{(n)}_s = \sum_{i<n}\gamma^i\delta_{s+i}$ with the
next-value rule of \cref{eq:gae}, and $M_n = \max(0, A^{(1)}_s, \dots, A^{(n)}_s)$. The
reference code's clipped running sum equals $A^{(n)}_s - M_n$, hence
\begin{equation}
  \hat G^{(n)}_s = M_n - A^{(n)}_s, \qquad
  \hat G_s = \sum_{n=1}^{N_s-1}(1-\lambda)\lambda^{n-1}\hat G^{(n)}_s
            + \lambda^{N_s-1}\hat G^{(N_s)}_s,
\end{equation}
with $N_s$ the number of steps from $s$ to the episode's end or the chunk end. This form
is exact at $\lambda = 1$, where the code's $1/(1-\lambda)$ form divides by zero.
% src: docs/paper/score-spec.md section 1.6; src/apexzero/audit/estimators.py (mna_terms)
The reference sums $\hat G$ over a fixed 512-step rollout; our windows vary in length, so we
normalise by the step count (\dev{11}), which ranks identically when $T$ is fixed. The
reference gate $\ind[\Rmax > 0]$ is almost always true under dense reward, so we gate
on ever completing a window (\dev{12}) and log the literal gate.

\paragraph{Baselines} Every endpoint is also computed for four reference scores: the
agent-only score $-\Jhat$, which knows nothing about the optimum; the policy-agnostic
scores $\Jstar$ and a geometric difficulty measure, which know nothing about the agent;
and a random score. An ``oracle-operational'' score $\Jstar - G_{\text{op}}$, the label
minus the return on the operational windows, is a \emph{reference, not a ceiling}:
critic-based scores pool earlier history and can beat it. The real ceiling on any score's
correlation with the noisy target is $\sqrt{\text{reliability}}$.
% src: docs/plan/two-track-plan.md section 2.2 ("Baselines")
\TODO{define the geometric difficulty measure}

\subsection{Two identities that split the error}
\label{sec:est-identities}

Two exact identities (experiment E2) turn RQ2's ``why'' into measured terms.
% src: docs/paper/score-spec.md sections 1.2 and 1.4; docs/plan/two-track-plan.md section 2.4 (E2)
Because $\max(a, 0) = (|a| + a)/2$, for every episode
\begin{equation}
  \mathrm{PVL}_e = \tfrac12\,\mathrm{L1}_e + \tfrac12\,\overline{A}_e,
  \label{eq:pvl-identity}
\end{equation}
so PVL differs from L1 only through the mean advantage, which is a measure of critic bias
on the level. For the unit-consistent MaxMC, with $\Jhat$ estimated on the disjoint
target windows,
\begin{equation}
  S_{\text{MaxMC-unit}} - \Reg
  = \underbrace{\big(\Rgmax - \Jstar\big)}_{\text{ceiling gap}}
  + \underbrace{\Big(\Jhat - \tfrac{1}{E}\textstyle\sum_e V_{t_0(e)}\Big)}_{\text{critic bias}}.
  \label{eq:maxmc-identity}
\end{equation}
The ceiling gap is at most zero unless the agent beats the label. We predict that on hard
levels it dominates the error of unit-consistent MaxMC (\cref{sec:prereg}). With staleness
(the age of a stored score) and success saturation (the tie fraction of $p(1-p)$), these
terms are the candidate explanations RQ2 tests.

\subsection{Aggregation and implementation checks}
\label{sec:est-aggregation}

In the fresh audit (E1), the \emph{operational} score of a level at a checkpoint uses
windows 1--4, which mimics what a curriculum sees; MaxMC's maximum is taken over those
four windows, which is exactly JaxUED's branch for a new level (\dev{23}).
\emph{High-budget} scores are cross-fitted: odd windows for the score, even windows for
the target. In training, a score is written when an episode ends, from the episodes of
one slot that end in the same chunk.
% src: docs/paper/score-spec.md section 1.8
Every score is computed from the same per-window records by one implementation, so the
in-training scorer and the offline audit cannot drift apart.
% src: src/apexzero/audit/estimators.py (module docstring)

The unit tests reproduce the spec's hand-computed vectors E-1--E-6, E-8--E-11, E-13--E-16
and E-19; PVL matches JaxUED's own function on completed episodes; and the closed-form
MNA equals a port of the reference code to $10^{-10}$ for $\lambda \in \{0.9, 0.95,
0.98\}$.
% src: tests/test_estimators.py (test names E1 ... E19; test_pvl_matches_jaxued_on_completed_episodes;
%      test_E15_mna_closed_form_equals_reference_port, atol 1e-10)
\TODO{implement spec tests E-7 (R\textsubscript{max} history), E-12 (MNA incomplete tail),
E-17 (chunk property test against JaxUED's accumulator), E-18 (online equals offline) and
the sampler tests P-1--P-9 and S-1--S-4; the samplers themselves are not built yet}

\subsection{Deviations from the reference code}
\label{sec:est-deviations}

\Cref{tab:deviations} is the complete list of departures from the pinned reference code,
reproduced from the scorer specification. Most follow from one fact: our episodes are
\num{1200}-step windows that span several 128-step training chunks, while every
reference scores whole episodes inside one rollout. The window-end row has no number
because it follows the reference.
% src: docs/paper/score-spec.md section 3 (table reproduced row by row)

{\footnotesize
\setlength{\tabcolsep}{4pt}
\renewcommand{\arraystretch}{1.15}
\begin{longtable}{@{}>{\raggedright\arraybackslash}p{1.05cm}>{\raggedright\arraybackslash}p{1.95cm}>{\raggedright\arraybackslash}p{3.2cm}>{\raggedright\arraybackslash}p{3.45cm}>{\raggedright\arraybackslash}p{2.6cm}@{}}
\caption{Deviations from the reference code (score spec, section 3). ``Ours'' describes
the racing audit and curricula. DEV-14 applies only if the ensemble heads share the main
critic's target; the current code trains each head on its own GAE target, as minimax does,
so DEV-14 no longer applies.}
\label{tab:deviations}\\
% src: docs/paper/score-spec.md section 3; DEV-14 status: src/apexzero/agents/ppo.py (gae_step,
%      "Ensemble heads get their own GAE targets"), docs/STATUS.md (K5 row)
\toprule
ID & Item & Reference & Ours & Reason \\
\midrule
\endfirsthead
\toprule
ID & Item & Reference & Ours & Reason \\
\midrule
\endhead
\midrule
\multicolumn{5}{r@{}}{\emph{continued on next page}}\\
\endfoot
\bottomrule
\endlastfoot
DEV-1 & Episode and rollout structure & each update resets envs to their levels and runs
  \code{num\_steps} with AutoReplay; scores cover the episodes completed in that rollout &
  asynchronous \num{1200}-step windows across 128-step chunks; per-slot carries; write at
  episode end & \num{1200}-step windows exceed the 128-step chunk \\
DEV-2 & $\lambda$-sum inside a scored episode & cut only at the rollout end, which affects
  only the dropped tail & cut and bootstrap at every chunk end; E1 emulates this with a
  phase offset $\varphi_e$; an unchunked variant is stored & the plan's protocol; Craftax
  precedent \\
DEV-3 & Parameter snapshots per episode & 1 (maze) & up to 11 & the policy updates between
  chunks; Craftax precedent \\
DEV-4 & Incomplete episodes & dropped (MNA includes them); $-\infty$ or 0 or NaN when none
  complete & carried until done; never scored at a chunk end & windows span chunks \\
DEV-5 & Episodes per score & all completed in the rollout ($\ge 1$ at maze defaults) & 1 per
  write in training (grouped within a chunk); 4 in E1 & window length \\
DEV-6 & Units & no reward normalisation; scores reuse the PPO's advantage tensor &
  $V := \sigma\hat V$; separate raw-unit scoring GAE & our PPO scales rewards
  (\code{ppo.py} L385--386) \\
-- & Window end (mode F) & time limit terminal & \textbf{same} (headline); mode B
  (bootstrap $V^f$) stored as a variant & follows the reference \\
DEV-8 & $\lambda$ & PPO $\lambda$ (maze 0.98) & PPO $\lambda = 0.95$ (+ variants) & same
  rule, different parameter \\
DEV-9 & Success & count of $r > 0$ steps, or GoalR & per-window indicator, S1 or S2 &
  dense reward \\
DEV-10 & Learnability trials & one AutoReplay stream, $\le 10$ episodes & $M$ independent
  windows & \num{1200}-step episodes \\
DEV-11 & MNA normalisation & $\Sigma$ over a fixed $T = 512$ steps & step-pooled mean
  (paper Eq.~12) & $T_e$ varies; identical ranking when $T$ is fixed \\
DEV-12 & MNA $\hat C$ & $\ind[\Rmax > 0]$ & $\ind[\text{ever S1}]$;
  $\hat C_{\text{lit}}$ logged & $\Rmax > 0$ is almost always true under dense reward \\
DEV-13 & Ensemble architecture & heads are the critic, with a shared trunk getting
  gradients; head 0 drives GAE & 5 extra heads on a stop-gradient trunk; main critic
  excluded; 64-unit ReLU & heads must not change the audited critic \\
DEV-14 & Ensemble targets (only if not fixed) & own GAE target per head & main critic's
  target (\code{ppo.py} L452--454) & current code; fix recommended \\
DEV-15 & Replay decision & one Bernoulli per update for the whole batch & one per slot per
  chunk & asynchronous resets \\
DEV-16 & $\perp$ mechanism & whole update discarded; no gradient step on DR updates &
  per-transition loss mask; a gradient step every chunk on $\sim$80\% of the data;
  normalisers masked & asynchronous mix inside each batch \\
DEV-17 & Weights during draws & re-evaluated per draw (sequential scan, 32 draws) & frozen
  per chunk ($\sim$220 draws); commits in $(t, j)$ order & buffer operations at chunk
  granularity \\
DEV-18 & Stale replay write & cannot happen (synchronous) & dropped if the slot's level was
  evicted & asynchrony \\
DEV-19 & Concurrent $\Rmax$ & stale read, last write wins (\code{update\_batch}) &
  sequential read of the latest value & asynchrony; the maximum is preserved \\
DEV-20 & Buffer and level source & 4000; generator; exact-field duplicate check & 1000;
  \num{4000}-level pool; duplicate check by \code{level\_id} & plan C \\
DEV-21 & SFL & $N$ fresh levels; $L$ steps; per-update reassignment; NaN not handled &
  $N$ = shuffled pool; $M$ windows; per-reset draw; $T$ in chunks; NaN asserted & finite
  pool; windows \\
DEV-22 & MaxMC-unit & -- & new estimator: discounted, $-\infty$ init, $V(o_0)$, no time
  average & the plan's variant \\
DEV-23 & E1 $\Rmax$ history & buffer history & windows 1--4 (= the new-level branch) &
  fresh audit \\
\end{longtable}
}
